\documentclass[9pt, mathserif, professionalfonts]{beamer}


\usepackage[ngerman]{babel}

\usepackage{fontspec}

\setsansfont{Tex Gyre Heros}


\setlength{\parindent}{0pt}
\usepackage[babel,german=quotes]{csquotes}
\usepackage{eurosym}

\usepackage{amsmath}
\usepackage{mathtools}
\usepackage{amssymb}
\usepackage{xfrac}
\usepackage{unicode-math}
\setmathfont{Tex Gyre Pagella Math}

\usepackage[locale=DE, separate-uncertainty=true, per-mode=fraction]{siunitx}

\usepackage{booktabs}
\usepackage{tabulary}
\usepackage{threeparttable}

\usepackage[]{graphicx}
\usepackage[labelfont=bf]{caption}
\setbeamertemplate{caption}[numbered]

\usepackage[]{hyperref}
\usepackage{braket}
\usepackage{xcolor}
\usepackage{enumerate}

\usepackage{listings}
\usepackage{lstautogobble}

\lstnewenvironment{lstverbatim}[1][]{
  \lstset
  {
    autogobble,
    language={[LaTeX]TeX},                
    basicstyle=\ttfamily,
    texcsstyle=*\color{TUdarkgreen},
    keywordstyle=\color{TUorange},
    commentstyle=\color{blue},
    morekeywords = {hallo},
    otherkeywords={$,$, &, \{, \}, \[, \],  ^, _},
    moretexcs={text, num, SI, si, kilo, gram, befehl, subject, titlehead, publishers,
               metre, per, second, squared, micro, subtitle, maketitle, tableofcontents,
               subsection, subsubsection, part, chapter, paragraph, subparagraph,
               iint, iiint},
    columns=fullflexible,
    breaklines=true,
    keepspaces=true,
    aboveskip=-0.2em,
    belowskip=-0.2em
  }
}{}


\title{Verfassen wissenschaftlicher Texte mit \LaTeX}
\author{PeP et Al. toolbox}
\date{2014}
\institute[Pep et Al. e.V.]{\includegraphics[width=0.2\textwidth]{peplogo.pdf}}
\titlegraphic{./latex-vs-word.pdf}

\usetheme{TUDo}

\begin{document}

\begin{frame}
\setcounter{framenumber}{0}
    \titlepage
\end{frame}

\begin{frame}
    \frametitle{Übersicht}
    \tableofcontents
\end{frame}

\section{Grundlagen}
\begin{frame}
    \frametitle{Warum \LaTeX \ benutzen?}
    \begin{itemize}
        \item hervorragender Text- und Formelsatz
        \item automatisierte Erstellung von Inhalts- und Literaturverzeichnis
        \item Tex-Dateien sind reine Text-Dateien \\
              $\Rightarrow$ kleine Datein, gut für Versionskontrolle geeignet
        \item sehr gute Vorlagen für wissenschaftliche Arbeiten 
        \item aber auch: Notensatz, Präsentationen 
        \item kostenlos, OpenSource, ausgezeichnete Dokumentionen
        \item erweiterbar durch zahlreiche und mächtige Pakete
        \item aus allen geläufigen Betriebssystem verfügbar
        \item Ausgabe nach PDF
    \end{itemize}
\end{frame}

\begin{frame}[fragile]
    \frametitle{Das Dokument}
    Diese drei Zeilen braucht jedes \LaTeX-Dokument:
    \begin{columns}[T]
        \begin{column}{0.4\textwidth}
            \begin{block}{\LaTeX-Code}
                \begin{lstverbatim}
                \documentclass{scrartcl}
                    % Präambel
                    .
                    .
                \begin{document}
                    % Inhalt des Dokuments
                    .
                    .
                \end{document}
                \end{lstverbatim}
            \end{block}
        \end{column}
        \begin{column}{0.55\textwidth}
            \begin{block}{documentclass}
                Hier wird die Dokumentenklasse definiert. \\
                Es folgt die \alert{Präambel} des Dokuments. 
                Hier werden globale Optionen gesetzt und zusätzliche Pakete eingebunden.\\
            \end{block}

            \vspace{0.55cm}

            \begin{block}{document-Umgebung}
                Hier wird das eigentliche Dokument erstellt.
            \end{block}
        \end{column}
    \end{columns}
\end{frame}

\begin{frame}[fragile]
    \frametitle{Syntax: Befehle}
    \LaTeX-Befehle beginnen stets mit einem Back-Slash.

    Obligatorische Argumente stehen in \{ \}, optionale Argumente stehen in [ ].
    \begin{columns}[T]
        \begin{column}{0.47\textwidth}
            \begin{block}{\LaTeX-Code}
                \begin{lstverbatim}
                \befehl[optional]{obligatorisch}
                \documentclass[paper=a4]
                              {scrartcl}
                \frac{1}{2}
                \end{lstverbatim}
            \end{block}
        \end{column}
        \begin{column}{0.47\textwidth}
            \begin{block}{Erklärung}
                Beispiel \\
                Setzt die Dokumentenklasse auf \emph{scrartcl} und das Papierformat auf DIN\,A4. \\
                Es gibt auch Befehle mit zwei oder mehr Pflichtargumenten, z.B. der Bruch.
            \end{block}
        \end{column}
    \end{columns}
\end{frame}

\begin{frame}[fragile]
    \frametitle{Standardpakete}
    Die hier aufgezählten Pakete sollten immer geladen werden, da sie wesentliche Funktionen bieten und wichtige Einstellungen vornehmen.
    \begin{columns}[T]
        \begin{column}{0.47\textwidth}
            \begin{block}{Paket}
                \begin{lstverbatim}
                \usepackage{polyglossia}
                \setmainlanguage{german}
                \usepackage{fontspec}



                \usepackage[unicode, pdfusetitle]{hyperref}
                \end{lstverbatim}
            \end{block}
        \end{column}
        \begin{column}{0.47\textwidth}
            \begin{block}{Funktion}
                Deutsche Spracheinstellungen für das Dokument. \\
                Schriftpaket. Wichtig, damit Umlaute korrekt dargestellt werden.
                Erlaubt einfaches ändern der verschiedenen Schriftarten.\\
                Für Hyperlinks im Dokument (z.B. Inhaltsverzeichnis $\rightarrow$ Kapitel).
            \end{block}
        \end{column}
    \end{columns}
\end{frame}

\begin{frame}[fragile]
    \frametitle{KOMA-Script-Klassen \hfill \doc{http://ftp.fernuni-hagen.de/ftp-dir/pub/mirrors/www.ctan.org/macros/latex/contrib/koma-script/doc/scrguide.pdf}{KOMA-Skript}}
    Stellt die wichtigen Klassen \texttt{scrartcl}, \texttt{scrreprt} und \texttt{scrbook} zur Verfügung. Sehr gute Vorlagen schon mit den Standardeinstellungen, schnell global mit Klassenoptionen anpassbar.
    \begin{block}{fürs Praktikum empfohlenene Einstellungen}
        \begin{lstverbatim}
        \documentclass[
            parskip=half,          % Absätze durch halbe Leerzeile
            captions=tableheading, % Spacing für Tabellenüberschriften
            headsepline,           % Linie zwischen Kopfzeile und Text
            titlepage=firstiscover % Titleseite ist Deckblatt
            ]{scrartcl}
        \end{lstverbatim}
    \end{block}
\end{frame}


\section{Text erstellen}

\begin{frame}[fragile]{Text schreiben}
  \begin{block}{Beispiel}
    \begin{lstlisting}
      % Präambel
      \begin{document}
        Hallo, Welt!

        Dies ist ein dummer Beispieltext.
        Er soll zeigen, dass \LaTeX sich nicht um
        Zeilenumbrüche im Code    oder    zuviele
        Leerzeichen kümmert.

        Ein Absatz wird mit einer leeren Code-Zeile
        markiert.
      \end{document}
    \end{lstlisting}
  \end{block}
\end{frame}

\begin{frame}[fragile]{Konventionen für Text}
  \begin{itemize}
    \item Höchstens ein Satz pro Code-Zeile
    \item Absätze werden durch eine Leerzeile markiert
    \item Im Fließtext sollten keine Umbrüche mit \verb+\\+ erzwungen werden
  \end{itemize}
  \begin{alertblock}{Sonderzeichen}
    Viele Sonderzeichen sind \LaTeX-Steuerzeichen.
    Damit diese im Text genutzt werden können, muss meist ein \verb+\+ vorangestellt oder ein Befehl genutzt werden.
  \end{alertblock}
  \begin{CodeExample}{0.6}
    \begin{lstlisting}
      \% \& \_ \textbackslash \$ \{ \}
    \end{lstlisting}
  \CodeResult
    \% \& \_ \textbackslash\ \$ \{ \}
  \end{CodeExample}
\end{frame}

\begin{frame}[fragile]{Textauszeichnung}
  Änderungen der Schrifteigenschaften sind mit diesen Befehlen möglich:
  \begin{CodeExample}{0.60}
    \begin{lstlisting}
      \textit{kursiv} \emph{kursiv}
      \textbf{fett}
      \textbf{\textit{fett-kursiv}}
      \textrm{Serifen-Schrift}
      \texttt{Mono-Schrift}
      \textsf{Sans-Serif-Schrift}
      \textsc{Kapitälchen}
    \end{lstlisting}
  \CodeResult
    \textit{kursiv} \emph{kursiv} \\
    \textbf{fett} \\
    \textbf{\textit{fett-kursiv}} \\
    \textrm{Serifen-Schrift} \\
    \texttt{Mono-Schrift} \\
    \textsf{Sans-Serif-Schrift} \\
    \textsc{Kapitälchen}
  \end{CodeExample}

  \vspace{1em}
  Diese Befehle sollten sehr selten benutzt werden, semantischer Markup ist besser.
\end{frame}

\begin{frame}[fragile]{Schriftgrößen}
  Gelten immer für den aktuellen Block, z.\,B. in einer Umgebung oder zwischen \lstinline+{ }+
  \begin{CodeExample}{0.48}
    \begin{lstlisting}
      {\tiny tiny}
      {\small small}
      {\normalsize normal}
      {\large large}
      {\huge huge}
    \end{lstlisting}
  \CodeResult
    {\tiny tiny}
    {\small small}
    {\normalsize normal}
    {\large large}
    {\huge huge}
  \end{CodeExample}
  \vspace{1em}
  \begin{block}{Alle Größen}
    \begin{lstlisting}
      \tiny, \scriptsize, \footnotesize, \small, \normalsize, \large, \Large, \LARGE, \huge, \Huge
    \end{lstlisting}
  \end{block}
  Auch diese Befehle sollten nur über semantischen Markup benutzt werden.
\end{frame}

\begin{frame}[fragile]{Inhalt auslagern}
  \begin{block}{Code}
    \begin{lstlisting}
      \PassOptionsToPackage{
  german,
  unicode,
  pdfusetitle,
  colorlinks,
  linkcolor=vertexDarkRed,
  urlcolor=vertexDarkRed,
  citecolor=vertexDarkRed,
}{hyperref}
\PassOptionsToPackage{
  aux,
}{rerunfilecheck}

\documentclass[
  9pt,
  aspectratio=1610,
]{beamer}
\usetheme{vertex}

\usefonttheme{professionalfonts}
\usepackage[english, french, ngerman]{babel}

\usepackage{fontspec}
\setmainfont{Latin Modern Roman}[
  SmallCapsFont = {Latin Modern Roman Caps},
]
\newfontfamily{\latinmodernroman}{Latin Modern Roman}[
  SmallCapsFont = {Latin Modern Roman Caps},
]
\DeclareTextFontCommand{\latinmodernrm}{\latinmodernroman}
\newfontfamily{\latinmodernsans}{Latin Modern Sans}
\DeclareTextFontCommand{\latinmodernsf}{\latinmodernsans}
\newfontfamily{\latinmodernmono}{Latin Modern Mono}
\DeclareTextFontCommand{\latinmoderntt}{\latinmodernmono}

\usepackage{fontawesome5}

\usepackage{amsmath}
\usepackage{mathtools}
\usepackage{amssymb}
\usepackage{mleftright}

\usepackage[
  math-style=ISO,
  bold-style=ISO,
  sans-style=italic,
  nabla=upright,
  partial=upright,
  mathrm=sym,
]{unicode-math}
\setmathfont{Latin Modern Math}[version=normal]
\setmathfont{XITS Math}[version=xits]
\setmathfont{XITS Math}[version=xitsss1, StylisticSet=1]
\mathversion{normal}

\usepackage[
  locale=DE,
  separate-uncertainty=true,
  per-mode=symbol-or-fraction,
]{siunitx}

\usepackage[
  version=4,
  math-greek=default,
  text-greek=default,
]{mhchem}

\usepackage[
  german=quotes,
  autostyle,
]{csquotes}
\usepackage{xfrac}

\usepackage{tabularray}
\UseTblrLibrary{booktabs, siunitx, varwidth}
\usepackage{threeparttable}

\usepackage{graphicx}

\usepackage[
  compatibility=false,
]{caption}
\usepackage{subcaption}

\usepackage{xcolor}
\usepackage{metalogo}
\usepackage{pdflscape}

\usepackage{fancyvrb}
\usepackage[highlightmode=immediate]{minted}
\setminted{
  autogobble,
  breaklines,
  stripnl=true,
}
\usemintedstyle{toolbox}

\usepackage[
  theorems,
  many,
]{tcolorbox}

\usepackage{tikz}
\usetikzlibrary{
  arrows,
  arrows.meta,
  graphs,
  graphdrawing,
  positioning,
  shadows,
  shapes,
}

\usegdlibrary{trees}
\usepackage[compat=1.1.0]{tikz-feynman}

\usepackage[
  shortcuts,
]{extdash}

\usepackage[noframe]{showframe}
\usepackage{bookmark}

\tcbset{
  noparskip,
  colback=white,
  colframe=vertexDarkGrey,
  boxsep=1mm,
  coltext=black,
  coltitle=white,
  left=0mm,
  right=0mm,
  top=0.5mm,
  toptitle=-1mm,
  bottom=0mm,
  bottomtitle=-1mm,
  sharp corners,
  before upper={\raggedright},
  before title={\raggedright\lstset{texcsstyle=*\color{white}}},
}

\lstset{
  autogobble,
  language={[LaTeX]TeX},
  basicstyle=\ttfamily,
  texcsstyle=*\color{vertexDarkRed},
  keywordstyle=\color{cyan},
  commentstyle=\color{darkgray},
  otherkeywords={$,$, &, \{, \}, \[, \],  ^, _},
  morekeywords={}, % do not remove even though empty
  columns=fullflexible,
  breaklines=true,
  keepspaces=true,
  showlines=true,
  aboveskip=0em,
  belowskip=0em,
}

\input{settings-lst-cs.tex}
\input{settings-lst-literate.tex}

% http://tex.stackexchange.com/a/267701/32665
% "Rehook" literate char table.
\makeatletter
\lst@AddToHook{SelectCharTable}
    {\ifx\lst@literate\@empty\else
         \expandafter\lst@Literate\lst@literate{}\relax\z@
     \fi}
\makeatother

\tikzset{
  invisible/.style={opacity=0,text opacity=0},
  visible on/.style={alt={#1{}{invisible}}},
  alt/.code args={<#1>#2#3}{%
    \alt<#1>{\pgfkeysalso{#2}}{\pgfkeysalso{#3}} % \pgfkeysalso doesn't change the path
  },
}
\setbeamertemplate{bibliography item}{\insertbiblabel}

\ExplSyntaxOn

\RenewDocumentEnvironment {block} {m o} {
  \IfValueTF{#2}{
    \begin{tcolorbox}[
      adjusted~title=#1,
      #2,
    ]
  }{
    \begin{tcolorbox}[
      adjusted~title=#1,
    ]
  }
}{
  \end{tcolorbox}
}

\RenewDocumentEnvironment {alertblock} {m o} {
  \IfValueTF{#2}{
    \begin{tcolorbox}[
      adjusted~title=#1,
      colframe=vertexDarkRed,
      #2,
    ]
  }{
    \begin{tcolorbox}[
      adjusted~title=#1,
      colframe=vertexDarkRed,
    ]
  }
}{
  \end{tcolorbox}
}

\RenewDocumentEnvironment {exampleblock} {m o} {
  \IfValueTF{#2}{
    \begin{tcolorbox}[
      adjusted~title=#1,
      colframe=blue!80!black,
      #2,
    ]
  }{
    \begin{tcolorbox}[
      adjusted~title=#1,
      colframe=blue!80!black,
    ]
  }
}{
  \end{tcolorbox}
}

\newcounter{heightgroup}

\NewDocumentEnvironment {CodeExplanation} {o m O{Code} O{Erklärung} D<>{}} {
  \NewDocumentCommand \Explanation {D<>{}} {
      \end{tcolorbox}
    \end{column}
    \begin{column}{
      \IfValueTF {#1} {
        #1
      }{
        \fp_eval:n {1 - 0.02 - #2}
      }
      \textwidth
    }
    \begin{tcolorbox}[
      equal~height~group=\theheightgroup,
      adjusted~title=#4,
      valign=top,
      ##1,
    ]
  }
  \begin{columns}[onlytextwidth, t]
    \begin{column}{#2 \textwidth}
      \begin{tcolorbox}[
        adjusted~title=#3,
        equal~height~group=\theheightgroup,
        valign=top,
        #5,
      ]
}{
      \end{tcolorbox}
    \end{column}
  \end{columns}
  \stepcounter{heightgroup}
}

\NewDocumentEnvironment {CodeExample} {o m O{Code} O{Ergebnis} D<>{}} {
  \NewDocumentCommand \CodeResult {D<>{}} {
      \end{tcolorbox}
    \end{column}
    \begin{column}{
      \IfValueTF {#1} {
        #1
      }{
        \fp_eval:n {1 - 0.02 - #2}
      }
      \textwidth
    }
      \begin{tcolorbox}[
        adjusted~title=#4,
        colframe=gray!40,
        coltitle=black,
        equal~height~group=\theheightgroup,
        valign=top,
        ##1,
      ]
        \begin{EmulateArticle}
  }
  \begin{columns}[onlytextwidth, t]
    \begin{column}{#2 \textwidth}
      \begin{tcolorbox}[
        adjusted~title=#3,
        equal~height~group=\theheightgroup,
        valign=top,
        #5,
      ]
}{
        \end{EmulateArticle}
      \end{tcolorbox}
    \end{column}
  \end{columns}
  \stepcounter{heightgroup}
}

\NewDocumentEnvironment {Packages} {} {
  \begin{block}{Benötigte~Pakete}
}{
  \end{block}
}

\makeatletter
\NewDocumentEnvironment {EmulateArticle} {} {
  \rmfamily
  \let\textrm=\latinmodernrm
  \let\textsf=\latinmodernsf
  \let\texttt=\latinmoderntt
  \renewcommand\thempfootnote{\arabic{mpfootnote}}

  \color{black}
  \setbeamercolor{item}{fg=black}
  \setbeamercolor{itemize/enumerate~body}{fg=black}
  \setbeamercolor{itemize/enumerate~subbody}{fg=black}
  \setbeamercolor{itemize/enumerate~subsubbody}{fg=black}
  \setbeamercolor{description~item}{fg=black}
  \setbeamercolor{enumerate~item}{fg=black}
  \setbeamercolor{itemize~item}{fg=black}
  \setbeamercolor{normal~text}{fg=black}
  \setbeamercolor{block~body}{fg=black}
  \setbeamerfont{item}{family=\rmfamily, size=\normalsize}
  \setbeamerfont{itemize/enumerate~body}{family=\rmfamily, size=\normalsize}
  \setbeamerfont{itemize/enumerate~subbody}{family=\rmfamily, size=\normalsize}
  \setbeamerfont{itemize/enumerate~subsubbody}{family=\rmfamily, size=\normalsize}
  \setbeamerfont{description~item}{series=\bfseries}
  \setbeamertemplate{itemize~item}{\bullet}
  \setbeamertemplate{itemize~subitem}{--}
  \setbeamertemplate{itemize~subsubitem}{\textasteriskcentered}
  \setbeamertemplate{enumerate~item}{\theenumi.}
  \setbeamertemplate{enumerate~subitem}{\alph{enumii})}
  \setbeamertemplate{enumerate~subsubitem}{\roman{enumiii}.}

  \setbeamerfont{footnote}{family=\rmfamily}
  \setbeamerfont{footnote~mark}{family=\rmfamily}

  \setbeamertemplate{caption}[numbered]
  \setbeamercolor{caption}{fg=black}
  \setbeamerfont{caption}{family=\rmfamily}
  \setbeamercolor{caption~name}{fg=black}
  \setbeamerfont{caption~name}{family=\rmfamily, series=\bfseries}

  % from amsmath.dtx
  % \def\maketag@@@#1{\hbox{\m@th\normalfont#1}}
  \def\maketag@@@##1{\hbox{\m@th\normalfont\rmfamily##1}}
  \hypersetup{linkcolor=black}
}{
}
\makeatother

\NewDocumentCommand \lstMakeSettingsInline {}
{
  language={[gnu]make},
  basicstyle=\ttfamily\color{black},
  otherkeywords={:, |, \\},
}

\NewDocumentCommand \lstMakeSettings {}
{
  \expandafter\lstset\expandafter{\lstMakeSettingsInline}
}

\NewDocumentCommand \lstmakeinline {O{}}
{
  \expandafter\lstinline\expandafter[\lstMakeSettingsInline,#1]
}

\box_new:N \l_listings_boxed_box
\lstnewenvironment{lstmake}[1][]{
  \lstMakeSettings
  \lstset{#1}
  \hbox_set:Nw \l_listings_boxed_box
}{
  \hbox_set_end:
  \leavevmode\box_use:N \l_listings_boxed_box
}

\NewDocumentCommand \lstPythonSettingsInline {}
{
  language=python,
  basicstyle=\ttfamily\color{black},
  otherkeywords={True, False, None, as},
  showstringspaces=false,
  stringstyle=\color{vertexDarkRed},
}

\NewDocumentCommand \lstPythonSettings {}
{
  \expandafter\lstset\expandafter{\lstPythonSettingsInline}
}

\NewDocumentCommand \lstinputpython {O{} m}
{
  \lstPythonSettings
  \hbox_set:Nw \l_listings_boxed_box
  \lstinputlisting[#1]{#2}
  \hbox_set_end:
  \leavevmode\box_use:N \l_listings_boxed_box
}

\NewDocumentCommand \lstpythoninline {O{}}
{
  \expandafter\lstinline\expandafter[\lstPythonSettingsInline,#1]
}

\lstnewenvironment{lstpython}[1][]{
  \lstPythonSettings
  \lstset{#1}
  \hbox_set:Nw \l_listings_boxed_box
}{
  \hbox_set_end:
  \leavevmode\box_use:N \l_listings_boxed_box
}

\NewDocumentCommand \lstMatplotlibrcSettings {}
{
  \lstset{
    basicstyle={\ttfamily\color{black}},
    morecomment={[l]\#},
    otherkeywords={:, \{, \}},
  }
}

\NewDocumentEnvironment {CenterStrip} {O{\textwidth} m}
{
  \begin{minipage}[c][#2\baselineskip][c]{#1}
}{
  \end{minipage}
}

\NewDocumentCommand \OverfullCenter {+m}
{
	\noindent
	\makebox[\linewidth]{#1}
}

\tikzstyle{buttonstyle} = [
  align=center,
  rectangle,
  fill=black!10,
  draw=black!80,
  drop~shadow,
  font={\bfseries},
]

\NewDocumentCommand \button {m}
{
  \tikz[
    baseline={([yshift=-.8ex]current~bounding~box.center)},
  ]{
    \node[buttonstyle] {\normalsize #1};
  }
}

\NewDocumentCommand \doc {m m}
{
  \button{\href{#1}{Doku:~\texttt{#2}}}
}

% used to show headline at beginning of each logical section
\NewDocumentCommand \headlineframe {m}
{
  \begin{frame}
    \begin{center}
      \Huge\color{vertexDarkRed}#1
    \end{center}
  \end{frame}
}

\NewDocumentCommand \removedisplayskip {}
{
  \setlength\abovedisplayskip{0ex}
  \setlength\belowdisplayskip{0ex}
  \setlength\abovedisplayshortskip{0ex}
  \setlength\belowdisplayshortskip{0ex}
}

\makeatletter
\seteverylogo
{
  \setlogodrop{0.7ex}
  \setLaTeXa{{\scshape a}}
  \fp_compare:nNnTF {\f@size} > {10}
  {
    \setlogokern{La}{-0.33em}
  }
  {
    \fp_compare:nNnTF {\f@size} < {10}
    {
      \setlogokern{La}{-0.4em}
    }
    {
      \setlogokern{La}{-0.36em}
    }
  }
}
\makeatother

\let\TeXold=\TeX
\def\TeX{\latinmodernrm{\TeXold}}
\def\eTeX{$ε$-\TeX}
\let\LaTeXold=\LaTeX
\def\LaTeX{\latinmodernrm{\LaTeXold}}
\let\XeTeXold=\XeTeX
\def\XeTeX{\latinmodernrm{\XeTeXold}}
\let\XeLaTeXold=\XeLaTeX
\def\XeLaTeX{\latinmodernrm{\XeLaTeXold}}
\let\LuaTeXold=\LuaTeX
\def\LuaTeX{\latinmodernrm{\LuaTeXold}}
\let\LuaLaTeXold=\LuaLaTeX
\def\LuaLaTeX{\latinmodernrm{\LuaLaTeXold}}
\def\pdfTeX{\latinmodernrm{pdf\TeX}}
\def\BibTeX{\latinmodernrm{\textsc{Bib}\TeX}}
\def\BibLaTeX{\latinmodernrm{Bib\LaTeX}}

% macro examples
\NewDocumentCommand \dif {m}
{
  \mathinner{\symup{d} #1}
}
\NewDocumentCommand \Dif {m}
{
  \mathinner{\symup{D} #1}
}
\NewDocumentCommand \del {m}
{
  \mathinner{\symup{\delta} #1}
}
\NewDocumentCommand \Del {m}
{
  \mathinner{\symup{\Delta} #1}
}

\let\vaccent=\v
\RenewDocumentCommand \v {}
{
  \TextOrMath{
    \vaccent
  }{
    \symbf
  }
}

\let\ltext=\l
\RenewDocumentCommand \l {}
{
  \TextOrMath { \ltext }{ \mleft }
  % \mathopen{}\mathclose\bgroup\left
}
\let\raccent=\r
\RenewDocumentCommand \r {}
{
  \TextOrMath { \raccent }{ \mright }
  % \aftergroup\egroup\right
}

\AtBeginDocument{
  \let\symRe=\Re
  \RenewDocumentCommand \Re {}
  {
    \operatorname{Re}
  }
  \let\symIm=\Im
  \RenewDocumentCommand \Im {}
  {
    \operatorname{Im}
  }
}

\DeclarePairedDelimiter{\abs}{\lvert}{\rvert}
\DeclarePairedDelimiter{\norm}{\lVert}{\rVert}
\DeclarePairedDelimiter{\bra}{\langle}{\rvert}
\DeclarePairedDelimiter{\ket}{\lvert}{\rangle}
\DeclarePairedDelimiterX{\braket}[2]{\langle}{\rangle}{#1\delimsize|#2}

\ExplSyntaxOff


\author{PeP et al.\ Toolbox Workshop}
\date{2017}
\institute[Pep et al.\ e.V.]{\includegraphics[width=0.4\textwidth]{logos/pep.pdf}}

      \begin{document}
        \input{Teil1.tex}
        \input{Teil2.tex}
        % .
      \end{document}
    \end{lstlisting}
  \end{block}
  \begin{itemize}
    \item Verschachtelung möglich
    \item Zur Aufteilung größerer Dokumente (z.B. diese Präsentation)
    \item Für häufig wiederverwendeten \LaTeX-Code (header, Erläuterungen zu Fehlerrechnung, …)
    \item Für per Skript erzeugte Tabelleninhalte
  \end{itemize}
\end{frame}

\begin{frame}[fragile]{Anführungszeichen}
  Die richtigen Anführungszeichen, wo die Satzzeichen hingehören und vieles mehr hängt von der Sprache ab.
  So macht man es richtig:
  \begin{Packages}
    \begin{lstlisting}
      \usepackage[autostyle]{csquotes}    % nach polyglossia
      \setotherlanguages{english, french} % andere Sprachen laden.
    \end{lstlisting}
  \end{Packages}
  \begin{CodeExample}{0.60}
    \begin{lstlisting}
      foo \enquote{bar} baz
      \enquote{foo \enquote{bar} baz}
      \textenglish{\enquote{foo}}
      \textfrench{\enquote{foo}}
      \textcquote{root}{foo}
    \end{lstlisting}
  \CodeResult
    foo \enquote{bar} baz \\
    \enquote{foo \enquote{bar} baz} \\
    \textenglish{\enquote{foo}} \\
    \textfrench{ \enquote{foo}} \\
    \textcquote{root}{foo}
    % work around polyglossia bug #68
    \directlua{polyglossia.desactivate_frpt()}
  \end{CodeExample}
\end{frame}

\section{Zahlen und Einheiten}

\begin{frame}[fragile]{Zahlen und Einheiten}
  \begin{itemize}
    \item Regeln zur Benutzung der SI-Einheiten: \\
      \url{http://www.bipm.org/utils/common/pdf/si_brochure_8_en.pdf}
    \item Einheiten werden aufrecht gesetzt
    \item Zwischen Zahl und Einheit steht ein kleines Leerzeichen
    \item Ab 5 Stellen wird ein kleines Leerzeichen als 1000er Trennzeichen genutzt:
  \end{itemize}

  \begin{CodeExample}{0.7}[Zahl mit Einheit]
    \begin{lstlisting}
      $5\,\mathrm{kg}$
    \end{lstlisting}
    \CodeResult
    $5\,\mathrm{kg}$
  \end{CodeExample}
  \begin{CodeExample}{0.7}[Zahl mit mehr als vier Stellen]
    \begin{lstlisting}
      $10\,000$
    \end{lstlisting}
    \CodeResult
    $10\,000$
  \end{CodeExample}
  \begin{CodeExample}{0.7}[Zehnerpotenz mit Unsicherheit in Klammern]
    \begin{lstlisting}
      $(5{,}34 \pm 0{,}54) \cdot 10^{-3}\,\mathrm{GeV}$
    \end{lstlisting}
    \CodeResult
    $(5{,}34 \pm 0{,}54) \cdot 10^{-3}\,\mathrm{GeV}$
  \end{CodeExample}
  
  \onslide<2->{%
  \begin{center}
    \LARGE  Das muss einfacher gehen
  \end{center}
  }
\end{frame}

\begin{frame}[fragile]{
  Das \texttt{siunitx}-Paket
  \hfill
  \doc{http://mirrors.ctan.org/macros/latex/contrib/siunitx/siunitx.pdf}{siunitx}
}
\begin{itemize}
    \item \texttt{siunitx} stellt Befehle zur Verfügung, die das korrekte Setzen von Zahlen und Einheiten stark vereinfachen
    \item Funktioniert in Fließtext und Matheumgebung
    \item[$\color{vertexDarkRed}\Rightarrow$] Dieses Paket sollte \alert{immer} und für \alert{jede} Zahl mit oder ohne Einheit verwendet werden.
\end{itemize}
  \begin{Packages}
    \begin{lstlisting}
      \usepackage[
        locale=DE,
        separate-uncertainty=true,  % Immer Fehler mit ±
        per-mode=symbol-or-fraction, % m/s im Text, sonst \frac
        % alternativ:
        % per-mode=reciprocal,      % m s^{-1}
        % output-decimal-marker=.,         % . statt , für Dezimalzahlen
      ]{siunitx}
    \end{lstlisting}
  \end{Packages}
\end{frame}

\begin{frame}[fragile]{\texttt{siunitx}: Zahlen mit \lstinline+\\num+}
  \begin{CodeExample}{0.7}[Zahlen mit automatischen 3er-Gruppen]
    \begin{lstlisting}
      \num{1.23456}
      \num{987654321}
    \end{lstlisting}
  \CodeResult
    \strut
    \num{1.23456} \\
    \num{987654321}
  \end{CodeExample}
  \begin{CodeExample}{0.7}[Einfaches Eingeben von 10er Potenzen]
    \begin{lstlisting}
      \num{6.022e23}
    \end{lstlisting}
  \CodeResult
    \strut
    \num{6.022e23}
  \end{CodeExample}
  \begin{CodeExample}{0.7}[Angabe von Fehlern]
    \begin{lstlisting}
      \num{1.54 +- 0.1}
      \num{1.54(10)}
      \num{1.54 \pm 0.1}
      \num[separate-uncertainty=false]{1.54 +- 0.1}
      \num{3.5(1)e6}
    \end{lstlisting}
  \CodeResult
    \strut
    \num{1.54 +- 0.1} \\
    \num{1.54(10)} \\
    \num{1.54 \pm 0.1} \\
    \num[separate-uncertainty=false]{1.54 +- 0.1} \\
    \num{3.5(1)e6}
  \end{CodeExample}
\end{frame}

\begin{frame}[fragile]{\texttt{siunitx}: Einheiten mit \lstinline+\\si+}
  \begin{CodeExample}{0.8}[Einheiten]
    \begin{lstlisting}
      \si{\meter\per\second}
      \si[per-mode=fraction]{\meter\per\second}
      \si{\meter\per\second\squared}
      \si[per-mode=reciprocal]{\gram\per\cubic\centi\meter}
      \si{\kelvin\tothe{4}}
    \end{lstlisting}
  \CodeResult
    \strut
    \si{\meter\per\second} \\
    \si[per-mode=fraction]{\meter\per\second} \\
    \si{\meter\per\second\squared} \\
    \si[per-mode=reciprocal]{\gram\per\cubic\centi\meter} \\
    \si{\kelvin\tothe{4}}
  \end{CodeExample}
  \begin{CodeExample}{0.80}[\texttt{per-mode=symbol-or-fraction}]
    \begin{lstlisting}
      \begin{equation}
        \si{\kilo\gram\meter\per\second\squared}
      \end{equation}
      $\si{\kilo\gram\meter\per\second\squared}$
    \end{lstlisting}
  \CodeResult
    \removedisplayskip
    \begin{minipage}[c][3\baselineskip][c]{\textwidth}
      \begin{equation}
        \si{\kilo\gram\meter\per\second\squared}
      \end{equation}
    \end{minipage}
    $\si{\kilo\gram\meter\per\second\squared}$
  \end{CodeExample}
  \begin{CodeExample}{0.80}[Meter mal Sekunde oder Millisekunde?]
    \begin{lstlisting}
      \si{\milli\second}
      \si{\meter\second}
      \si[inter-unit-product=\cdot]{\meter\second}
    \end{lstlisting}
  \CodeResult
    \strut
    \si{\milli\second} \\
    \si{\meter\second} \\
    \si[inter-unit-product=\cdot]{\meter\second}
  \end{CodeExample}
\end{frame}

\begin{frame}[fragile]{\texttt{siunitx}: Zahl mit Einheit: \lstinline+\\SI+}
  \begin{CodeExample}{0.6}[\lstinline+\\SI+ {=} Kombination aus \lstinline+\\num+ und \lstinline+\\si+]
    \begin{lstlisting}
      \SI{5}{\percent}
      \SI{10}{\celsius}
      \SI{2.5(1)e6}{\kilo\gram\square\meter\per\second\squared}
    \end{lstlisting}
  \CodeResult
    \strut
    \SI{5}{\percent} \\
    \SI{10}{\celsius} \\
    \SI{2.5(1)e6}{\kilo\gram\square\meter\per\second\squared}
  \end{CodeExample}
  \begin{description}
    \item[1. Argument] Kann alles, was \lstinline+\num+ kann
    \item[2. Argument] Kann alles, was \lstinline+\si+ kann
  \end{description}
  \begin{CodeExample}{0.6}[Winkel]
    \begin{lstlisting}
      \ang{5;55;59}
    \end{lstlisting}
  \CodeResult
    \strut
    \ang{5;55;59}
  \end{CodeExample}
\end{frame}

\section{Formelsatz}
\subsection{im Fließtext}
\begin{frame}[fragile]
   \frametitle{\$ \dots \$-Umgebung}
   Für mathematische Symbole, Variablen und kleine Formeln im Fließtext. 

   \begin{block}{benötigte Pakete}
    \begin{lstverbatim}
        \usepackage{amsmath}
        \usepackage{mathtools}
        \usepackage{amssymb}
    \end{lstverbatim}
   \end{block}
    \begin{columns}[T]
        \begin{column}{0.6\textwidth}
            \begin{block}{\LaTeX-Code}
                \begin{lstverbatim}
                $x$ 
                $x^i$
                $x^12$ bzw. $x^{12}$ % Vorsicht
                $x_\text{max}$
                $U(t) = U_0 \cdot \cos(\omega t)$
                \end{lstverbatim}
            \end{block}
        \end{column}
        \begin{column}{0.35\textwidth}
            \begin{block}{Ergebnis}
                $x$  \\
                $x^i$ \\
                $x^12$ bzw. $x^{12}$ \\ 
                $x_\text{max}$ \\
                $U(t) = U_0 \cdot \cos(\omega t)$
            \end{block}
        \end{column}
    \end{columns}
\end{frame}

\input{./Inhalt/gleitumgebungen.tex}

\end{document}
